% (C) 2026 Ralf Stephan, in collaboration with Claude Code.
% Released under CC0 1.0 Universal (public-domain dedication).
% See https://creativecommons.org/publicdomain/zero/1.0/

\documentclass[11pt]{article}
\usepackage[T1]{fontenc}
\usepackage[utf8]{inputenc}
\usepackage{lmodern}
\usepackage[margin=2.5cm]{geometry}
\usepackage{amsmath,amssymb}
\usepackage{booktabs,tabularx}
\usepackage{braket}
\usepackage[hidelinks]{hyperref}

% Lean names and paths: verbatim, breakable at . _ /
\DeclareUrlCommand\lean{\urlstyle{tt}}
\def\UrlBreaks{\do\.\do\_\do\/}
\newcolumntype{L}{>{\raggedright\arraybackslash}X}
\newcolumntype{B}{>{\raggedright\arraybackslash}p{0.42\textwidth}}
\emergencystretch=6em

\title{Lean formalization of an Assumptions of Physics technical brief:\\ log}
\author{}
\date{}

\begin{document}
\maketitle

\noindent
This log records which statements of the technical brief in this directory were checked in Lean,
where the formal version differs from the brief, and any problems found in the brief.
The brief is available at \url{https://assumptionsofphysics.org/resources/briefs}.

\section*{Setup}

\begin{itemize}
\item Lean \texttt{v4.36.0-rc2}; Mathlib at commit \texttt{1b1ccf4b}; physlib (\texttt{master})
  and TauCeti (commit \texttt{e7e7f52}, on the same Mathlib) as further dependencies.
\item Build: \texttt{lake build} in the project root.
\item Every file is checked with \texttt{\#print axioms} on its main results. A result is
  recorded here only if it depends on no axioms beyond \lean{propext},
  \lean{Classical.choice} and \lean{Quot.sound}, and so has no \lean{sorry}.
\item Imports are kept to the Mathlib modules actually needed.
\item Material that belongs in Mathlib is kept in \lean{ForMathlib/}, in Mathlib's directory
  layout (library \lean{ForMathlib}), and imported by the brief files.
\end{itemize}


%======================================================================

\section{Equivalence of sup-norm, \texorpdfstring{$L^1$}{L1}, and \texorpdfstring{$L^2$}{L2}
  seminorms on the Schwartz space}

\paragraph{Files.}\leavevmode
\begin{itemize}
\item Files: \lean{briefs/SchwartzNorms/SchwartzNorms.lean},\newline
  \lean{briefs/SchwartzNorms/SchwartzNorms/GeneralLp.lean} (the Remark ``General $L^p$''),\newline
  \lean{briefs/SchwartzNorms/SchwartzNorms/FiniteMoments.lean} (the converse of the remark after
  eq.~(10))
\item Library: \lean{SchwartzNorms}
\item Imports: \lean{SchwartzMomentProblem} (for the moments in eq.~(10)) and three new
  \lean{ForMathlib} files:\newline
  \lean{ForMathlib/Analysis/LocallyConvex/WithSeminorms.lean} (bounds between seminorm
  families compose),\newline
  \lean{ForMathlib/Analysis/SpecialFunctions/ImproperIntegrals.lean}
  ($\int (1+x^2)^{-2} = \pi/2$, $\int_{|x|>1} x^{-2} = 2$),\newline
  \lean{ForMathlib/MeasureTheory/Integral/IntegralEqImproper.lean}
  ($\|g(x)\| \le \|g'\|_{L^1}$ for integrable $g$, $g'$);\newline
  \lean{GeneralLp.lean} imports only \lean{SchwartzNorms};\newline
  \lean{FiniteMoments.lean} imports \lean{SchwartzNorms} and four TauCeti modules:
  \lean{TauCeti.Analysis.Distribution.TemperedDistribution.Lp},
  \lean{TauCeti.Analysis.Distribution.TemperedDistribution.SmulLeft},
  \lean{TauCeti.Analysis.Sobolev.WeakDeriv.Interval} and
  \lean{TauCeti.Analysis.Sobolev.WeakDeriv.TemperedDistribution}
\end{itemize}

\paragraph{Setting.} For $f : \mathbb{R} \to \mathbb{C}$ and $g = x^\alpha \partial^\beta f$
(\lean{weighted}), the seminorms (1)--(3) are \lean{supNorm}, \lean{l1Norm} and \lean{l2Norm},
with values in $[0, \infty]$ (\lean{eLpNorm} for $L^1$ and $L^2$). Smoothness is
\lean{ContDiff} of order $\infty$. On $\mathcal{S}(\mathbb{R})$, the $L^1$ and $L^2$ seminorms
are real seminorms \lean{l1Seminorm}, \lean{l2Seminorm}, collected in the families
\lean{l1Family}, \lean{l2Family}; the sup family is Mathlib's \lean{schwartzSeminormFamily}.

\paragraph{Results.}
\begin{center}
\begin{tabularx}{\textwidth}{@{}BL@{}}
\toprule
Brief & Lean \\
\midrule
eq.~(6), sup controls $L^2$ & \lean{l2Norm_sq_le} \\
eq.~(7), $L^2$ controls $L^1$ & \lean{l1Norm_le} \\
eq.~(9), $L^1$ controls sup & \lean{supNorm_le} (with $\|f\|^{L^1}_{\alpha,\beta} < \infty$);
  \lean{not_supNorm_le_const} (the hypothesis is needed) \\
Theorem 5, (i) $\Leftrightarrow$ (ii) $\Leftrightarrow$ (iii) & \lean{tfae_finite} \\
definition (4) is Mathlib's $\mathcal{S}(\mathbb{R})$ & \lean{exists_schwartzMap_iff},
  \lean{supNorm_eq_seminorm} \\
Step 4, same topology & \lean{withSeminorms_l1_l2}, from \lean{isBounded_sup_l2},
  \lean{isBounded_l2_l1}, \lean{isBounded_l1_sup} and \lean{WithSeminorms.congr_of_cycle} \\
eq.~(10) & \lean{enorm_moment_le} (with the factor $|\hbar|^\beta$) \\
Remark ``General $L^p$'': the $L^p$ seminorms, $1 \le p \le \infty$, generate the same topology & \lean{withSeminorms_lp}, from \lean{isBounded_sup_lp} (\lean{eLpNorm_weighted_le}), \lean{isBounded_lp_l1} (\lean{eLpNorm_weighted_one_le}) and \lean{isBounded_l1_sup}; \lean{lpFamily_one}, \lean{lpFamily_two}, \lean{lpFamily_top} (\lean{eLpNorm_top_weighted}) for $p = 1, 2, \infty$ \\
remark after eq.~(10), converse: finite moments imply $f \in \mathcal{S}(\mathbb{R})$ &
  \lean{exists_schwartzMap_iff_forall_memL2} (tempered distributions), from
  \lean{exists_schwartzMap_of_forall_memL2} and \lean{forall_memL2_of_schwartzMap};
  \lean{exists_schwartzMap_ae_eq} ($L^2$ functions); \lean{exists_schwartzMap_of_l2Norm_lt_top}
  (smooth functions) \\
\bottomrule
\end{tabularx}
\end{center}

\paragraph{General $L^p$ seminorms.} The Remark ``General $L^p$'' is formalized in
\lean{GeneralLp.lean}. For $1 \le p \le \infty$ the seminorms
$\|f\|^{L^p}_{\alpha,\beta} = \|x^\alpha \partial^\beta f\|_{L^p}$ (\lean{eLpNorm}, with the
essential supremum for $p = \infty$) form the family \lean{lpFamily}~$p$, and every such family
generates the topology of $\mathcal{S}(\mathbb{R})$ (\lean{withSeminorms_lp}). The cycle has
three steps.
\begin{itemize}
\item Sup controls $L^p$. With the weight $w_k(x) = (1+|x|)^{-k}$ (\lean{weight}) and
  $(1+|x|)^k \le 2^k(1+|x|^k)$,
  \[
    \|f\|^{L^p}_{\alpha,\beta} \le 2^k \|w_k\|_{L^p}
      \bigl(\|f\|^{\sup}_{\alpha,\beta} + \|f\|^{\sup}_{\alpha+k,\beta}\bigr)
  \]
  (\lean{eLpNorm_weighted_le}), and $\|w_k\|_{L^p} < \infty$ for $k$ large
  (Mathlib's \lean{Measure.HasTemperateGrowth.exists_eLpNorm_lt_top}).
\item $L^p$ controls $L^1$. Write $x^\alpha \partial^\beta f = w_k \cdot (1+|x|)^k x^\alpha
  \partial^\beta f$ and apply Hölder's inequality with the conjugate exponent $q$,
  $1/p + 1/q = 1$:
  \[
    \|f\|^{L^1}_{\alpha,\beta} \le 2^k \|w_k\|_{L^q}
      \bigl(\|f\|^{L^p}_{\alpha,\beta} + \|f\|^{L^p}_{\alpha+k,\beta}\bigr)
  \]
  (\lean{eLpNorm_weighted_one_le}), with $k$ chosen for $q$.
\item $L^1$ controls sup: Step~3 of the brief (\lean{isBounded_l1_sup}).
\end{itemize}
For $p = 1, 2$ the family is \lean{l1Family}, \lean{l2Family} by definition. For $p = \infty$ the
essential supremum of the continuous function $x^\alpha \partial^\beta f$ is its supremum
(\lean{eLpNorm_top_weighted}), so the family is the family of sup seminorms~(1)
(\lean{lpFamily_top}).

\paragraph{Finite moments imply Schwartz.} The remark after eq.~(10) says that finiteness of
all expectation values is equivalent to membership in $\mathcal{S}(\mathbb{R})$, and the brief
proves only the bound (10). The converse is formalized in \lean{FiniteMoments.lean}. For it to
say something, $X^\alpha$ and $\partial^\beta$ must act on states outside
$\mathcal{S}(\mathbb{R})$, so the state is an arbitrary tempered distribution $u$, differentiated
in the sense of distributions (Mathlib's \lean{TemperedDistribution.derivCLM} and
\lean{TemperedDistribution.smulLeftCLM}). Then
\[
  u \in \mathcal{S}(\mathbb{R}) \iff x^\alpha \partial^\beta u \in L^2 \text{ for all }
  \alpha, \beta
\]
(\lean{exists_schwartzMap_iff_forall_memL2}). The finiteness conditions are the $L^2$ seminorms
(3), $\|x^\alpha \partial^\beta u\|_{L^2}^2 = \braket{X^\alpha \partial^\beta u | X^\alpha
\partial^\beta u}$. For Schwartz states and $\hbar \ne 0$ these are, up to the factor
$\hbar^{2\beta}$, the expectation values of the symmetric products $P^\beta X^{2\alpha} P^\beta$
(\lean{SchwartzMomentProblem.matrixElement_P_pow_X_pow_P_pow}).
For an $L^2$ state $f$ the conclusion is that $f$ agrees almost everywhere with a Schwartz
function (\lean{exists_schwartzMap_ae_eq}). For smooth $f$ it is Theorem~5 together with
\lean{exists_schwartzMap_iff} (\lean{exists_schwartzMap_of_l2Norm_lt_top}).

The proof of $\Leftarrow$ has three steps.
\begin{itemize}
\item The derivatives $\partial^\beta u$ are $L^2$ functions $g_\beta$, so $g_{\beta+1}$ is the
  weak derivative of $g_\beta$ (TauCeti's
  \lean{hasWeakLineDerivOn_of_lineDerivOp_toTemperedDistribution_eq}), and $g_\beta$ agrees
  almost everywhere with $G_\beta(x) = c_\beta + \int_0^x g_{\beta+1}$ (TauCeti's
  \lean{HasWeakLineDerivOn.exists_ae_eq_add_intervalIntegral_top}).
\item The $G_\beta$ are continuous and $G_\beta(x) = c_\beta + \int_0^x G_{\beta+1}$, so
  $G_\beta' = G_{\beta+1}$ everywhere by the fundamental theorem of calculus. Hence $G_0$ is
  smooth with $G_0^{(\beta)} = G_\beta$.
\item The $L^2$ function given for $x^\alpha \partial^\beta u$ agrees almost everywhere with
  $x^\alpha G_0^{(\beta)}$ (TauCeti's \lean{Lp.ae_eq_of_toTemperedDistribution_apply_eq}), so
  all $L^2$ seminorms of $G_0$ are finite and $G_0 \in \mathcal{S}(\mathbb{R})$ by Theorem~5.
\end{itemize}
The direction $\Rightarrow$ holds because derivatives and multiplication by $x^\alpha$ of the
distribution of a Schwartz function are those of the function
(\lean{forall_memL2_of_schwartzMap}).

\paragraph{Differences from the brief.}
\begin{itemize}
\item The seminorms take values in $[0, \infty]$, so that Theorem~5 can be stated for every
  smooth $f$. On $\mathcal{S}(\mathbb{R})$ they are finite and give real seminorms.
\item ``Generate the same locally convex topology'' is stated as \lean{WithSeminorms}: each
  family induces Mathlib's topology on $\mathcal{S}(\mathbb{R})$. Step~4's ``this relation is
  transitive'' is \lean{Seminorm.IsBounded.comp}, new in \lean{ForMathlib}.
\item In Step~3, $g(x) = \int_{-\infty}^x g'$ uses Mathlib's improper FTC and the fact that an
  integrable $g$ with integrable $g'$ tends to $0$ at $-\infty$.
\item The converse in the remark after eq.~(10) (finite expectation values imply membership in
  $\mathcal{S}(\mathbb{R})$) needs a meaning for $X^\alpha P^\beta f$ outside
  $\mathcal{S}(\mathbb{R})$. The formalization takes the state to be a tempered distribution and
  the condition to be $x^\alpha \partial^\beta u \in L^2$, i.e.\ the finiteness of the $L^2$
  seminorms (3), rather than of the expectation values $\braket{f|X^\alpha P^\beta|f}$ of
  eq.~(10). The brief does not say which orderings of $X$ and $P$ count.
\item The Remark ``General $L^p$'' says that ``the same arguments extend''. The formal proof
  replaces the $L^2$ step by Hölder's inequality and the weight $1/(1+x^2)$ by
  $(1+|x|)^{-k}$, with $k$ depending on $p$. It goes $\sup \to L^p \to L^1 \to \sup$, so it
  reuses Step~3 but not Steps~1 and~2.
\end{itemize}

\paragraph{Problems found in the brief.}
\begin{itemize}
\item Eq.~(9) needs $\|f\|^{L^1}_{\alpha,\beta} < \infty$, so that $g \to 0$ at $-\infty$. The
  brief has it from (ii), but (9) is stated as an inequality between seminorms, and it is false
  without it: for $f = 1$, $\alpha = \beta = 0$, the left side is $1$ and the right side is $0$.
\item In eq.~(7) the coefficient of $\|f\|^{L^2}_{\alpha+1,\beta}$ is $2$; the argument given
  ($\sqrt{a+b+c} \le \sqrt a + \sqrt b + \sqrt c$ with $b = 2(\|f\|^{L^2}_{\alpha+1,\beta})^2$)
  gives $\sqrt 2$. The stated bound is still true, and is the one formalized.
\item Eq.~(10) omits the factor $|\hbar|^\beta$ from $P^\beta = (-i\hbar\partial)^\beta$.
\end{itemize}

\end{document}
